\section{Construction of the operational partitions}
\label{sec:partitions}

We now prove Proposition~\ref{tab:partition-input}.  At a component
creation or change, the new cells must be short and padded, must contain
the inherited accounting classes, and must not put surviving points from
two different cells of one old child into a common new cell.  The danger
trains of Construction~\ref{sp:danger-trains} allow us to satisfy these
requirements simultaneously.  We first bound the cost of keeping target
points away from too much inherited danger history, and then choose new
boundaries outside the inherited class hulls.

Throughout this section, every comparison is between centres of the same
exact width $r$.  We work chronologically, conditional on the entire
construction before a component creation or change at height $Y_P$.
All old partitions, weights, classes and histories are therefore fixed.
We assume the partition assertions already hold for the old children.
We retain the stronger induction that the actual arc length of an
occupied old cell, at any surviving centre-$j$ point, is at most
$2r\exp(-2E_gB_j)<1/4$.  The two directional boundary tests below
give this bound at each creation; later projections and deletions
do not enlarge a cell.  At the first creation there are no old cells.
The sparse estimates in Proposition~\ref{sp:sparse-estimates} apply to this
prior history, and Construction~\ref{sp:danger-trains} supplies its
deterministic trains.  No choice made below changes an earlier partition
or an earlier coarsening edge.

\subsection{Weighted target-grid counts}

Let $M=L_j$ be a target centre.  At the current projection write
\[
 W'=M/c,
 \qquad
 R_{\mathrm{done}}(W')
   =\prod_{\substack{p\mid W'\\p\ \mathrm{processed}}}R_p,
\]
where $c$ is the processed part of one of its retained good row masks.
Fix the remaining mask and hence its phase.  Its current grid can be written as
$x_z=(z+\xi)/W'$, with $z$ running through $W'$ consecutive integers.
Denote the actual raw weight there by $w_z$, setting $w_z=0$ at a dead
entry.  The remaining-mask probability is not included in $w_z$.

A source train generated by centre $L=L_i$ in an earlier sparse bin $Y'$
has denominator $V'=L/d$.  Here $d$ includes the single omission $q$
at its generating step, and $q$ is the largest prime divisor of $d$.
For one of its prescribed offsets $y_0$, its centres are
$(n+y_0)/V'$, $n\in\mathbb Z$.  As usual, put
\[
 G=(L,M),\qquad U=L/G,\qquad T=M/G,
 \qquad m=\frac{W'}{(W',V')}.
\]
In particular $m\mid Td$.  Counts involving the two grids use periodic
lifts, so a source centre is counted every time it lies in the indicated
interval, even if several such centres are near the same target point.

\begin{lemma}[A target grid against a source train]
\label{part:target-grid}
For every sufficiently small fixed $\epsilon>0$, every fixed constant
$C_0$, and $0<R\le C_0$, one has
\begin{equation}
\begin{split}
 \frac{1}{W'}\sum_{z\bmod W'}w_z
 \sum_{n\in\mathbb Z}
 \mathbf 1_{\{|x_z-(n+y_0)/V'|\le C_0rR\}}
 \ll_{\epsilon,C_0}
 R_{\mathrm{done}}(W')e^{CY'}
 \left(V'rR+\frac{(2UT)^\epsilon}{m}\right).
\end{split}
\label{eq:target-grid}
\end{equation}
The constant $C$ in the harmonic factor is absolute.  The estimate is
uniform in the preceding history, fixed masks and phases.  It holds also
at a new insertion, where only static target fields are present.
\end{lemma}

\begin{proof}
Use the single-row product upper fields from
Corollary~\ref{tab:intermediate-fields}.  First fix the static anchor choices, select
one active birth exclusion at each applicable prime, and work inside
their average.  Soft tests, holes and intervening discards can be
dropped.  The resulting product upper version dominates the actual
weight after the static average is taken.  Each retained field at a
processed prime divisor of $W'$ has its prescribed one-row mean $R_p$.

Set $P_0=\rho\log(2UT)$, where $\rho>0$ will be sufficiently small in
terms of $\epsilon$.  The proximity multiplicity on the left of
\eqref{eq:target-grid} is a query modulo $m$ in $z$.  Drop the fields
at processed primes $p\mid m$ with $p>P_0$.  Relative to their ordinary
means, this costs at most
\[
 \prod_{\substack{p\mid m\\p>P_0}}R_p^{-1}
 \le \exp\bigl(CY'+O_\rho(1)\bigr).
\]
Indeed $m\mid Td$, primes of $d$ lie below the end of bin $Y'$, and
their reciprocal sum is $O(Y')$.  For primes of $T$ above
$\rho\log(2UT)$, the bounds on their size and on
$\sum_{p\mid T}\log p$ give a bounded reciprocal sum, with constant
depending only on $\rho$.

Keep every multiplier as the function defined by the actual prior
history; removing a field from an upper product does not recompute
any matching plan.  Average the retained large fields in decreasing
prime order.  At an adaptive full factor $p^\nu$, the current grid
still contains that full power.  Since $p\nmid m$, $p^\nu\mid V'$
also, so translation by $1/p^\nu$ preserves both the grid and the
proximity query.  Every retained earlier field is invariant on this
orbit, and the adaptive field supplies exactly its mean $R_p$.
An adaptive omission supplies only the scalar $R_p$ when its residual
power is positive, and supplies no field when that power is zero.
At a retained static prime of residual exponent $\mu>0$, translation
by $1/p^\mu$ fixes the query and all other retained static residues,
while its numerator step is a unit modulo $p$; hence its single
exclusion again averages to $R_p$.  All static primes of this target
precede its adaptive range, so no unaveraged adaptive field remains
when this static averaging is performed.  No independence of the
matching plans is used.

Deleting any of the high fields has left the low fields themselves
unchanged.  If $p_*\le P_0$ is the last adaptive update at such a
prime, let $Q_*$ be its historical post-update core.  The low fields
are individually invariant under translation by $1/Q_*$.  Intersect
this group with the current grid: the intersection consists of
translations by multiples of $1/(W',Q_*)$.  Thus their joint numerator
period divides
\[
 D=\frac{W'}{(W',Q_*)}\le\frac{M}{Q_*}
   \le\exp(O(p_*))\le\exp(O(P_0)).
\]
Here $Q_*\mid M$ and $W'\mid M$ give the first inequality prime by
prime, and \eqref{eq:core-loss} gives the second.  If no low adaptive
update exists, only static residue fields and scalars remain; one can
take $D=\prod_{p\mid W',\ p\le P_0}p$, which divides $W'$ and has
the same bound.  Hence the period estimate remains valid after
projection and after arbitrary deletion of the specified high fields.
Choose $\rho$ small in terms of $\epsilon$ to obtain
$D\ll_\epsilon(2UT)^\epsilon$.  Bounded $P_0$, including $U=T=1$,
is harmless; with no nonconstant low fields take $D=1$.

Write $g=(W',V')$.  In units of the source grid, the target coordinate
has the form
\[
 V'x_z-y_0
   =\frac{V'/g}{m}z+\text{a fixed shift},
 \qquad (V'/g,m)=1.
\]
Conditional on a residue of $z$ modulo $D$, its residues modulo $m$
run uniformly through a coset whose mesh, after this unit
multiplication, is $(D,m)/m$.  Thus the average number of source
centres in the query is
\[
 O\left(V'rR+\frac{(D,m)}{m}\right),
\]
uniformly in that residue modulo $D$.  This is the ordinary uniform
grid count, including multiplicities.  Multiply by the low product
field and average over its residues.  Its product mean, together with
the already averaged large fields, supplies the ordinary factors in
$R_{\mathrm{done}}(W')$, apart from the bounded restoration cost
already displayed.  Finally average the static anchor choices.  This
proves \eqref{eq:target-grid}.
\end{proof}

The next mask calculation separates an unrestricted source sum from
the target assignment.  This distinction will be important when the
target density is too small to discard its assignment restriction.
For a processed mask $c$, let $c_{\le Y'}$ be its part supported on
primes up through the end of bin $Y'$.

\begin{lemma}[Source masks and target deficit currency]
\label{part:mask-sums}
For the source masks of bin $Y'$,
\[
 \sum_d\frac1d\le e^{CY'}.
\]
At every fixed target mask $c$, put
\[
 E_{\mathrm{ext}}=
 \frac{W'}{(W',L)}=\frac{T}{(T,c)}.
\]
Then
\begin{equation}
 \sum_d\frac1m
 \ll E_{\mathrm{ext}}^{-1}
       \tau(U)\tau(c_{\le Y'})e^{CY'}.
\label{part:conditional-mask-sum}
\end{equation}
Under the unrestricted deficit law on the processed target mask,
\begin{equation}
 \mathbb E\bigl[(T,c)\tau(c_{\le Y'})\bigr]
 \le \tau(T)^C e^{CY'}.
\label{part:target-mask-moment}
\end{equation}
All the harmonic exponents here have absolute coefficients.
\end{lemma}

\begin{proof}
We may extend each nonnegative source sum to all divisors of $L$
supported below the end of bin $Y'$.  The first inequality follows by
multiplying the geometric sums and using the reciprocal-prime bound.
The generating prime introduces no additional sum: it is the largest
prime divisor of $d$, so is determined by that mask.

For the second inequality, fix a prime and write
$n=\nu_p(L)$ and $w=\nu_p(W')$.  If the exponent of $p$ in $d$ is
$j$, then its exponent in $m$ is
\[
 (w-n)_++\bigl(j-(n-w)_+\bigr)_+.
\]
The first term is the exponent in $E_{\mathrm{ext}}$.  Writing
$a=(n-w)_+$, the local sum after removal of that term is at most
\[
 \sum_{j=0}^n p^{-(j-a)_+}
 \le \frac{a+1}{1-1/p}.
\]
Moreover
\[
 a\le \nu_p(U)+\nu_p(c),
 \qquad
 a+1\le (\nu_p(U)+1)(\nu_p(c)+1).
\]
Only primes in the indicated source segment contribute the second
factor involving $c$.  Multiplication proves
\eqref{part:conditional-mask-sum}.  The formula for
$E_{\mathrm{ext}}$ is also immediate prime by prime from $W'=M/c$.

To prove \eqref{part:target-mask-moment}, use independence of the
deficit coordinates.  At $p^t\parallel T$, with deficit exponent
$J$, its geometric tail gives
\[
 \mathbb E\bigl[p^{\min(t,J)}(J+1)\bigr]
 \ll (t+1)^2.
\]
For example, the terms with $J<t$ sum to at most a constant times
$\sum_{j<t}(j+1)$, and the geometric tail with $J\ge t$ costs
$O(t+1)$.  If the divisor factor is absent, this estimate only
improves.  At primes not dividing $T$, the needed divisor moment is
$1+O(1/p)$.  Their product over the source segment is $e^{CY'}$,
while the other factors are bounded by a fixed power of $\tau(T)$.
This proves the assertion.
\end{proof}

Let $b$ denote a target outer mask and $I$ its remaining part.  Its
baseline first-mass coefficient satisfies
\begin{equation}
 rW'R_{\mathrm{done}}(W')\pi_I(b)
       =k_M\pi_M(cb).
\label{part:target-currency}
\end{equation}
Summing this over the original assigned masks gives $m_M=k_M\delta_M$.
Consequently the main term in \eqref{eq:target-grid} may always retain
the target assignment.  In contrast, when we sum its error without
that restriction, \eqref{part:target-mask-moment} introduces the
relative factor $1/(T\delta_M)$, up to arbitrarily small powers of
$2UT$.  At a fixed assigned target mask we can instead use
\eqref{part:conditional-mask-sum}, with its factor
$E_{\mathrm{ext}}^{-1}\tau(c_{\le Y'})$, and incur no inverse
assignment density.

\subsection{History tests and their first-mass cost}

For an already generated source centre-bin history $(i,Y')$, put
\[
 \ell'=\exp\bigl(-.3\exp(\sigma_gY')\bigr),
 \qquad Z'=\exp\bigl(\exp(h_aY')\bigr).
\]
Use the complete deterministic supply from
Construction~\ref{sp:danger-trains}: for every permissible good
processed mask $d$ generated in this bin, take its list of at most
$Z'$ offsets, and put an interval of radius $r\ell'$ around each
point of the corresponding grid of denominator $L_i/d$.  Denote the
resulting periodic family of intervals by $\mathcal H_{i,Y'}$.
Completeness here means all permissible masks and prescribed offsets
of this already generated history, not future or unprocessed bins.
We do not restrict the supply to trains supporting edges that
happened to survive; later deletion of an entry never deletes a
part of this history.

For a target point $x$ and a specified relative scale $0<t\ll1$,
call a tested bin \emph{threatening at scale $t$} if
\[
 |\log\ell'|\le 100|\log t|.
\]
Apply the following test separately for each tested source
centre-bin pair $(i,Y')$:
\begin{enumerate}
 \item In a threatening bin, discard $x$ if it lies within
 $(\ell'+2t)r$ of any centre of a train in
 $\mathcal H_{i,Y'}$.
 \item In a nonthreatening bin, discard $x$ if the sum of the lengths
 of the intervals of $\mathcal H_{i,Y'}$ meeting
 $[x-2r,x+2r]$ exceeds $e^{-Y'}tr$.
\end{enumerate}
The second sum may count overlaps with multiplicity, which only
strengthens the test.  Both tests use periodic lifts.

At the present component change there are two kinds of tests.  Every
target centre $j$ uses its intrinsic scale
\[
 t_j=\exp(-2E_gB_j)
\]
against all its already generated histories in the component.  In
addition, a point belonging to an old child $D$ uses the child-gap
scale
\[
 t_D=\exp(-1.5E_2Y_D),
\]
where $Y_D$ is that child's creation or last change height, against
histories of \emph{different} old children.  In such a foreign-child
comparison, with $Y_{\mathrm{prev}}=Y_P/2$, the prior component
separation gives
\begin{equation}
 \log(UT)>\beta Y_{\mathrm{prev}},
 \qquad
 Y_{\mathrm{prev}}\ge\max(Y',Y_D),
 \qquad B_i,B_j\le Y_{\mathrm{prev}}/A_0.
\label{part:foreign-separation}
\end{equation}

\begin{lemma}[Cost of a history test]
\label{part:history-test-cost}
For either kind of test, the first mass discarded from the tested
target array on $M=L_j$ by one source centre-bin pair $(i,Y')$ is at
most $e^{-cY'}m_M$.  Here $c>0$ can be taken sufficiently large for
the centre-count summations below, after the fixed parameter choices
and then a sufficiently small $k_0$.  The bound is uniform in the
prior history and does not depend on the slot cap $N_0$.
\end{lemma}

\begin{proof}
\emph{Threatening main term.}\quad
In a threatening bin, apply \eqref{eq:target-grid} with
$R=\ell'+2t$, sum the at most $Z'$ offsets for each $d$, and then
sum $d$.  Keeping the target assignment in the main term and using
$V'=L/d$ and Lemma~\ref{part:mask-sums} gives the relative main cost
\begin{equation}
 \ll e^{CY'}Z'k_L(\ell'+2t).
\label{part:threat-main}
\end{equation}
The threatening inequality implies
$t\le\exp(-.003\exp(\sigma_gY'))$.  Also
$B_i\le Y'/A_0$ and $\log k_L\le .2B_i+O(1)$.
Since $h_a<\sigma_g$, expression
\eqref{part:threat-main} is smaller than any required fixed
exponential in $-Y'$ once the heights are large.

\emph{Nonthreatening tests.}\quad
For a nonthreatening bin, an interval of radius $r\ell'$ meeting
$[x-2r,x+2r]$ has its centre within $3r$ of $x$.
Apply \eqref{eq:target-grid} at this constant relative radius, and
use Markov's inequality for the length threshold $e^{-Y'}tr$.
The length per counted interval is $2r\ell'$.  The factor $e^{Y'}$
from the threshold is absorbed in the harmonic exponential.  In the
main term retain the target assignment; only in the error extend
the target-mask sum to the unrestricted law.  By
\eqref{part:target-mask-moment} and arbitrarily small divisor powers,
the relative cost is at most
\begin{equation}
 C_\epsilon\frac{\ell'}{t}e^{CY'}Z'
 \left(k_L+\frac{(2UT)^\epsilon}{T\delta_M}\right).
\label{part:nonthreat-cost}
\end{equation}
All small powers of $2UT$ here may be enlarged and renamed
$\epsilon$.

We check that \eqref{part:nonthreat-cost} is negligible for both
scales.  Equal widths imply
\begin{equation}
 \log U-\log T=\log k_L-\log k_M,
 \qquad |\log U-\log T|\ll B_i+B_j.
\label{part:equal-width-ratios}
\end{equation}
Thus, for $\epsilon<1/2$, the ratio
$(2UT)^\epsilon/T$ is bounded by an exponential in a small fixed
multiple of $B_i+B_j$.  By \eqref{eq:centre-lower}, the density loss is at most
$\delta_M^{-1}\le \exp(.32B_j^*)\le\exp(.64B_j)$.
Write $A=|\log\ell'|$ and $a=|\log t|$.  In the present case
$A>100a$, so $\ell'/t=\exp(-(A-a))$ has a fixed large margin over
$\exp(-a)$.  At the intrinsic scale, $a=2E_gB_j$; at the
child-gap scale, $a=1.5E_2Y_D$ and $B_j\le Y_D/A_0$.
In both cases this margin pays the displayed density and ratio
losses.  It also pays $e^{CY'}Z'$ and $k_L$, because
$\log Z'=\exp(h_aY')=o(\exp(\sigma_gY'))$.

\emph{Threatening error term.}\quad
It remains to estimate the threatening error, where the factor
$\ell'/t$ is unavailable.  Its unrestricted relative bound is
\begin{equation}
 C_\epsilon e^{CY'}Z'
       \frac{(2UT)^\epsilon}{T\delta_M}.
\label{part:threat-error}
\end{equation}
For a foreign-child test, the threatening inequality gives
$\exp(\sigma_gY')=O(Y_D)$ and hence
$\log Z'=o(Y_D)$.  Equations
\eqref{part:foreign-separation} and
\eqref{part:equal-width-ratios} show that the denominator $T$ pays
\eqref{part:threat-error}: its logarithmic saving is
$(1/2-\epsilon)\log(UT)$ up to $O(B_i+B_j)$, whereas the
remaining height costs are $CY'+o(Y_D)+O(B_j)$.
Taking $\beta$ sufficiently large absolutely absorbs these costs,
uniformly since $Y',Y_D\le Y_{\mathrm{prev}}$.

For an intrinsic test, the threatening inequality instead gives
\[
 Y'=O(1+\log B_j)=o(B_j),\qquad
 \log Z'=o(B_j),\qquad B_i/B_j=o(1).
\]
If $\log(UT)\ge C_0B_j$ for a sufficiently large fixed $C_0$,
the same use of \eqref{part:equal-width-ratios} makes
\eqref{part:threat-error} exponentially small in $B_j$.
Suppose therefore that $\log(UT)<C_0B_j$.  Retain the target
assignment and apply \eqref{part:conditional-mask-sum} separately
at each of its masks.  For every true assigned row $v$ at $M$,
equal widths rule out the width alternative in
\eqref{eq:exterior-gain}, so
\[
 \frac{v}{(v,L)}>\exp(.015B_j^*).
\]
Since $v\mid W'$, it follows that
\[
 E_{\mathrm{ext}}=\frac{W'}{(W',L)}
       \ge \frac{v}{(v,L)}\ge \exp(.01B_j).
\]
The reciprocal-prime sum in the initial segment through bin $Y'$
is $O(Y')=o(B_j)$.  The good-mask initial-segment condition
therefore gives $\tau(c_{\le Y'})\le e^{s'B_j}$.
Relative to this assigned mask's own baseline
\eqref{part:target-currency}, its error is consequently at most
\[
 C_\epsilon e^{CY'}Z'(2UT)^\epsilon
 E_{\mathrm{ext}}^{-1}\tau(c_{\le Y'})
 \le \exp\bigl(-.01B_j+s'B_j+\epsilon C_0B_j+o(B_j)\bigr).
\]
Here the factor $\tau(U)$ has been absorbed by enlarging the
arbitrarily small power.  Choose that power sufficiently small;
$s'$ was already chosen small absolutely.  This gives the required
saving without a factor $\delta_M^{-1}$.  Summing the retained
target assignments completes all cases.
\end{proof}

We next sum these estimates without charging an unchanged history
repeatedly.  At a fixed target centre, its intrinsic scale never
changes.  For each source centre-bin pair, perform the intrinsic
comparison at the first partition change after that history has
been generated and has joined the target's component.  A later
position is a descendant at exactly the same location, and the
tested family is the same full deterministic supply, so it cannot
newly fail that intrinsic test.  For foreign-child comparisons,
the charge is indexed by the fixed triple $(j,i,Y')$, not by a
child's changing name.  Components merge but never split.  If this
source history exists before the first join of $i$ and $j$, that
join is its only possible foreign comparison with the old target
child.  Afterward $i$ and $j$ belong to the same old child at every
later change.  If the history is generated after their join, it is
never foreign to $j$.  Thus changing child-gap scales introduce no
repeated charge.  No monotonicity of the foreign test as $t_D$
changes is being used.

At height $Y'$, the number of possible source centres is at most
$\exp(4Y'/A_0+O(1))$.  Summing
Lemma~\ref{part:history-test-cost} over sources and dyadic heights,
then over target first masses $m_M$, gives $o(W)$ as $k_0\to0$.
The same count, without target weights, bounds the total physical
length per unit period of all danger intervals by
\begin{equation}
 \sum_{i,Y'} e^{CY'}Z'k_{L_i}\ell'<.01,
\label{part:global-danger-length}
\end{equation}
after decreasing $k_0$ if necessary.  Indeed each source mask
contributes at most $2Z'(L_i/d)r\ell'$, and
$\sum_d1/d\le e^{CY'}$.  The double-exponential decay of $\ell'$
pays both this harmonic factor and the centre count.

All history tests are equivariant under the current core.  Each
source supply was invariant under its post-generation core, and
the current common core divides that earlier core.  The predicates
use only distances, lengths and the fixed centre scales.  The
preliminary survivors after these tests therefore retain the
required core symmetry.

\subsection{Inherited hulls and new boundaries}

We have bounded the mass needed to keep short segments around
surviving points mostly outside the inherited histories.  We now
use those free portions to place boundaries.  The construction is
performed separately for each exact width $r$ in the component.

For each nonempty inherited accounting class, order its preliminary
surviving positions within its occupied old cell, whose actual arc
length is less than $1/4$ by the stronger induction.  Define its
\emph{short hull} to be the closed arc between the first and last
of these positions.  Omit empty classes and use a singleton hull
when all positions coincide.

The preliminary positions, inherited hulls and preliminary gaps
between occupied old cells are determined before any jitter is
sampled; they are not recomputed after the subsequent pointwise
purges.  The following observation ensures
that avoiding these hulls is affordable.

\begin{lemma}[Location of inherited class hulls]
\label{part:inherited-hulls}
Every positive-length inherited class hull lies in the union of
the danger intervals belonging to its old child.  It lies in its
old cell, and the hull operation respects current-core translations.
\end{lemma}

\begin{proof}
An inherited equivalence can be witnessed by a historical chain
of actual coarsening edges.  The endpoints never change location
under subsequent projections.  Construction~\ref{sp:danger-trains}
covers each connecting segment by that child's stored history.
If intermediate labels have disappeared, retain their segments
in the witnessing chain: the full deterministic supply has not
been reduced.  The historical chains may backtrack, but their
images lie in one connected component $K$ of the old child's
finite closed danger union.  By
\eqref{part:global-danger-length}, $K$ is an arc of length less
than $.01$; a chain with nonzero winding would cover the circle
and is impossible.

The current class also lies in one occupied old cell $C$.  By the
stronger induction specified at the start of the section, its
actual arc length is at most $2r\exp(-2E_gB_j)<1/4$ at any
surviving centre-$j$ point.  Choose a cut outside $C\cup K$.
Both become real intervals, and the interval between the extreme
surviving class positions lies in both.  Its length is less than
$.01$, so it is the uniquely determined short circular hull.
Old padding keeps it inside the old cell, including with the
half-open endpoint convention.  All classes and histories are
equivariant, and this hull operation commutes with the current
core translations.  Singleton hulls require no historical segment
and satisfy the equivariance assertion directly.
\end{proof}

Write $Q$ for the core at the start of the present bin and put
$T_{\mathrm{core}}=1/Q$.  Subdivide one period into $N$ equal
microintervals, where
\[
 N=\left\lceil\frac{T_{\mathrm{core}}}{r\exp(-E_2Y_P)}\right\rceil,
 \qquad a_P=T_{\mathrm{core}}/N\asymp r\exp(-E_2Y_P).
\]
The stated comparability holds because $Q\mid M$ at every target
centre and hence
\[
 Qr\exp(-E_2Y_P)
 \le k_M\exp(-E_2Y_P)\ll1.
\]
Choose one independent uniform candidate boundary point, called a
\emph{raw line}, in each microinterval, and replicate the
configuration with period $1/Q$.  Skip a line if it
belongs to any closed inherited hull of the preliminary survivors.
The hulls used for this decision are now fixed; they are not
shrunk and the lines are not redrawn after later discards.
Figure~\ref{fig:inherited-hull} illustrates the distinction between
surviving labels and the hull protected by this rule.

\begin{figure}[htbp]
\centering
\setlength{\unitlength}{1pt}
\begin{picture}(390,150)
\put(0,136){\makebox(390,12)[l]{\small One historical connecting chain}}
\put(55,106){\line(1,0){230}}
\put(55,106){\circle*{5}}
\put(285,106){\circle*{5}}
\put(170,106){\makebox(0,0){$\times$}}
\put(55,86){\makebox(0,12){$x_-$}}
\put(285,86){\makebox(0,12){$x_+$}}
\put(170,117){\makebox(0,12){\small vanished connector}}
\put(0,66){\makebox(390,12)[l]{\small Preliminary survivors and fixed hull}}
\linethickness{1.4pt}
\put(55,42){\line(1,0){230}}
\linethickness{.4pt}
\put(55,42){\circle*{5}}
\put(285,42){\circle*{5}}
\put(145,27){\line(0,1){30}}
\put(345,27){\line(0,1){30}}
\put(145,7){\makebox(0,12){\small skipped line}}
\put(345,7){\makebox(0,12){\small not blocked here}}
\end{picture}
\caption{An inherited class viewed in a cut coordinate. A connector that
vanished before the preliminary tests no longer appears among the
surviving labels, but its connecting segments remain covered by the
stored train history. The surviving extremes therefore determine a
hull in which sampled raw lines are skipped. The outside candidate is
not blocked by this hull; other hulls are not shown. The preliminary
hull is fixed before the lines are sampled and is not recomputed after
later purges. Only one possible chain and its class are shown;
historical chains may backtrack and different classes' hulls may overlap.}
\label{fig:inherited-hull}
\end{figure}


The unskipped lines define the new cyclic partition.  If none
remain, use the single circle cell.  Use half-open cells with a
fixed orientation when assigning endpoints.  First discard every
entry within
\[
 \delta_P=r\exp(-1.5E_2Y_P)
\]
of a raw line, including a skipped one.  Conditional on the prior
history, each fixed entry has probability
$O(\delta_P/a_P)=O(\exp(-E_2Y_P/2))$ of this discard.

For each preliminary target point $x$ of centre $j$, require an
unskipped line in each directional segment of length $t_jr$
starting at $x$, and discard the entry on failure.  At least half
of either segment is available for such a line.  To see this,
Lemma~\ref{part:inherited-hulls} puts all the skipped hulls in the
danger histories.  The threatening tests exclude their intervals
from these segments: a train centre more than
$(\ell'+2t_j)r$ from $x$ cannot have its radius-$\ell'r$ interval
meet them.  For nonthreatening histories, the total covered length
is at most
\[
 t_jr\sum_{i,Y'}e^{-Y'}<\tfrac12t_jr.
\]
The source-centre count and the large minimum height justify the
last inequality.  Singleton hulls have zero length.

We record a quantitative non-repetition condition for both the
intrinsic segments and the child-gap segments used below.  Write
$\log k_M\le .2B_j+C_k$, where $C_k$ is an absolute constant.
Since $Qr\le k_M$, the choices of scales give
\[
 \begin{aligned}
  \frac{a_P}{T_{\mathrm{core}}}&\le
       \exp\bigl(C_k-(E_2-.2/A_0)Y_P\bigr),\\
  \frac{t_jr}{T_{\mathrm{core}}}&\le
       \exp\bigl(C_k-(2E_g-.2)B_j\bigr),\\
  \frac{t_Dr}{T_{\mathrm{core}}}&\le
       \exp\bigl(C_k-(1.5E_2-.2/A_0)Y_D\bigr)
       \quad\text{when }B_j\le Y_D/A_0.
 \end{aligned}
\]
Take the minimum bands and heights sufficiently large that each
right-hand side is at most $1/4$.  For either segment length
$s=t_jr$ or $s=t_Dr$, this ensures
$s+a_P\le T_{\mathrm{core}}/2<T_{\mathrm{core}}$.
Distinct periodic copies of one microinterval are separated by
a gap $T_{\mathrm{core}}-a_P>s$.  Thus the tested segment meets at most one
copy of each independently sampled microinterval, even when it
crosses the chosen fundamental-period boundary.  For any
measurable available subset $A$ of the segment, the probability
that no raw line falls in $A$ is
\[
 \prod_J\left(1-\frac{|A\cap J|}{a_P}\right)
 \le \exp\left(-\frac{|A|}{a_P}\right),
\]
where the product is over those microintervals.  Hence the
two-direction failure probability is at most
\begin{equation}
 2\exp\bigl(-c t_j\exp(E_2Y_P)\bigr).
\label{part:intrinsic-jitter}
\end{equation}
This estimate permits any fragmentation of the available set.
Since $B_j\le Y_P/A_0$ and $E_2\gg E_g$, it is uniformly
negligible.  Every occupied new cell now has actual arc length at
most $2t_jr<1/4$ as measured from each surviving point of centre $j$.
This proves the stronger occupied-cell induction used for the hulls.
In particular it has the required intrinsic span
$\ll r\exp(-E_gB_j)$.

At an insertion containing a high-type newborn, the entire
component is newborn.  Indeed the last dense bin agrees across
each component edge at that height: the edge ratio is too small
to change any exponent in that bin.  Every active centre in the
component consequently has the same last dense bin and starts
now.  There are thus no inherited hulls to skip.  With one raw
line in each microinterval, every cell has length at most
$2a_P\ll r\exp(-E_gY_P)$, which is the stronger span needed
for high-type direct births.

\subsection{Separation of old cells and completion of the induction}

It remains to ensure that one new cell cannot contain surviving
points from two different old cells of the same child.  Fix an
old child $D$ and its nonempty old cells, where nonemptiness is
measured using the preliminary survivors.  If there is only one
such cell there is nothing to prove.  Otherwise list these cells
cyclically.  Between consecutive ones, the arc from the last
preliminary surviving point of the first to the first of the
next has length at least $2t_Dr$.  Indeed the old padding places
both points at distance at least $t_Dr$ from the old boundaries
between them.

The child's own current hulls do not enter this gap, by
Lemma~\ref{part:inherited-hulls} and the choice of the surviving
extrema.  Consider the segment of length $t_Dr$ running from
either endpoint into the gap.  The foreign-child history tests
at that endpoint leave at least half its length outside all
other children's hulls, by exactly the coverage argument used
for the intrinsic segments.  The non-repetition condition proved
above applies here, since
\[
 Qt_Dr\le t_Dk_M\ll1,
 \qquad B_j\le Y_D/A_0
\]
for its endpoint centre $j$; in particular $t_Dr+a_P<T_{\mathrm{core}}$.

For each nonempty old cell test its two adjacent gaps in these
directions.  If either test finds no unskipped line, discard all
entries of this child in that old cell.  Always use the fixed
preliminary positions and gaps, even if other subsequent tests
have already removed some points.  The same independent-jitter
calculation gives a failure probability per old cell of at most
\begin{equation}
 2\exp\bigl(-c t_D\exp(E_2Y_P)\bigr).
\label{part:child-jitter}
\end{equation}
Since $Y_P\ge2Y_D$,
\[
 t_D\exp(E_2Y_P)
   =\exp\bigl(E_2Y_P-1.5E_2Y_D\bigr)
   \ge\exp(E_2Y_P/4).
\]
Thus this failure probability is uniformly negligible as well.

Two old cells of this child that still contain survivors now
have an unskipped line between their preliminary live positions
in each direction around the circle.  This remains true if
intermediate old cells have lost all their entries: the separating
lines were selected using the original gaps and are never
removed or redrawn.  Consequently surviving points from this
child in a common new cell come from one old cell.

Let $\mathcal F$ denote all data after the preliminary history tests
and before the current jitter draws.  The pointwise bounds for
raw-line padding and \eqref{part:intrinsic-jitter} multiply the
corresponding fixed preliminary weights.  For an old cell $C$, let
$M_C^0$ be its preliminary first mass, including the outer-mask
factors, and let $E_C$ be its gap-test failure event defined using
the fixed preliminary gaps and hulls.  Any earlier pointwise
discards only reduce the mass left in $C$, so its subsequent
cell-purge loss is at most $\mathbf1_{E_C}M_C^0$ for every
realization.  Consequently
\[
 \mathbb E[\text{cell-purge loss in }C\mid\mathcal F]
 \le \mathbb P(E_C\mid\mathcal F)M_C^0.
\]
We never condition on the outcomes of the other tests.  Their
correlation with $E_C$ is therefore irrelevant, and overlapping
loss estimates only overcount the actual discard.  The preliminary
old-cell label sets are disjoint, and their total first mass is at
most $W$.  Thus \eqref{part:child-jitter} can be summed with these
mass weights.  The total entering first mass at each height is
also at most $W$.  Summation over the finite dyadic chronology gives
$o(W)$, uniformly in its length.  Together with the history-test
loss this is the required negligible partition discard.

No unskipped boundary enters an inherited closed hull, so
restricting a class to survivors does not split it across new
cells.  Every surviving point has distance at least $\delta_P$
from all new boundaries.  At a later sparse step of this
unchanged component, of height $Y'\ge Y_P$, the relaxed
coarsening scale satisfies
\[
 r\exp\bigl(-\tfrac12\exp(\sigma_gY')\bigr)
   \ll r\exp(-1.5E_2Y_P)=\delta_P.
\]
Thus the new good coarsening edges cannot cross a boundary.
The inequalities hold from the chosen minimum height onward,
because $\exp(\sigma_gY_P)\gg E_2Y_P$.

Finally, all parts of the construction respect the core
symmetry.  Old objects remain equivariant for the smaller common
core, the hull and coverage operations commute with its
translations, and the realized raw-line configuration is copied
with exactly that period.  The directional and old-cell tests
are transported under the same translations.  This establishes
the equivariance of the realized partitions and discards, not
merely invariance of their probability law.

The new partitions contain inherited classes, separate surviving
old cells within each child, and have the required span, padding
and realized core equivariance.  We may now apply the prescribed
cap clipping inside these cells.  The old-child separation supplies the
clipping inequality from Section~\ref{sec:tables}; the clipping
loss is accounted for separately there.  Subsequent sparse
updates retain the cells, their padding and their inherited
classes as described above.  An unchanged component needs no
partition replacement and no fresh clipping.  The construction
therefore continues through every component creation and change,
and proves Proposition~\ref{tab:partition-input}.
